% ===========================================================================
% Appendix: Weak-field perturbative expansion of q^(t) - q^(tau)
% ===========================================================================
% Drop-in appendix section for the brachistochrone paper.
% Conventions: G = c = 1, mu = r_s/R = 2M/R.
% ===========================================================================

\appendix

\section{Weak-field perturbative expansion of the $t$--$\tau$ splitting}
\label{app:weakfield}

We derive an analytic expression for the difference
$\delta q \equiv q^{(t)} - q^{(\tau)}$ between the coordinate-time and
proper-time brachistochrone turning radii, accurate to second order in the
compactness parameter $\mu = r_s/R = 2M/R$.

\subsection{Setup and notation}

The TOV metric inside the star is
\begin{equation}
    ds^2 = -A(r)\,dt^2 + b(r)\,dr^2 + r^2\,d\Omega^2,
\end{equation}
with lapse $A(r) = e^{2\nu(r)}$ and radial metric
$b(r)=\bigl(1-2m(r)/r\bigr)^{-1}$.
We work with the dimensionless radial coordinate $x = r/R$, and define
\begin{align}
    \bar{m}(x) &= m(xR)/M, \\
    W(x) &= \int_x^1 \frac{\bar{m}(x')}{x'^2}\,dx',
    \label{eq:W_def}
\end{align}
so that $W(1)=0$ and $W(0) = \int_0^1 \bar{m}(x)/x^2\,dx > 0$.

\subsection{Weak-field expansion of the optical indices}

To leading order in $\mu$ the TOV equation for $\nu$ gives
\begin{equation}
    \nu(r) - \nu(R) = -\frac{\mu}{2}\,W(x) + O(\mu^2),
\end{equation}
so that $A(x) = A_0\,e^{-\mu W(x)}$ where $A_0 = A(R) = 1-\mu + O(\mu^2)$.
The surface lapse difference is
\begin{equation}
    A_0 - A(x) = A_0\bigl(1 - e^{-\mu W(x)}\bigr) \approx \mu W(x)\bigl(1 - \mu W(x)/2\bigr) + O(\mu^3).
\end{equation}
The radial metric factor expands as
\begin{equation}
    \sqrt{b(x;\mu)} = \frac{1}{\sqrt{1-\mu\bar{m}(x)/x}}
        \approx 1 + \frac{\mu\,\bar{m}(x)}{2x} + O(\mu^2).
    \label{eq:sqrtb}
\end{equation}

The two optical indices are
\begin{align}
    n_t(x;\mu) &= \frac{1}{\sqrt{A(A_0-A)}}, &
    n_\tau(x;\mu) &= \sqrt{\frac{A}{A_0-A}}.
\end{align}
Introducing the rescaled ``momentum''
$\tilde{p}^2(x;\mu) \equiv \mu\,n^2(x;\mu)\,x^2$ (finite for $\mu\to 0$),
one finds
\begin{align}
    \tilde{p}_t(x;\mu)   &= P_0(x)\Bigl[1 + \mu\Bigl(1 + \tfrac{3}{4}W(x)\Bigr)\Bigr] + O(\mu^2),
    \label{eq:pt_expand}\\
    \tilde{p}_\tau(x;\mu) &= P_0(x)\Bigl[1 - \tfrac{\mu}{4}\,W(x)\Bigr] + O(\mu^2),
    \label{eq:ptau_expand}
\end{align}
where
\begin{equation}
    P_0(x) = \frac{x}{\sqrt{W(x)}}
    \label{eq:P0}
\end{equation}
is the Newtonian ``momentum'' profile.

\subsection{The Fermat angle integrals and their perturbative expansion}

The turning radius $q^{(t)}R$ (or $q^{(\tau)}R$) is the root of the
brachistochrone condition
\begin{equation}
    \Phi^{t,\tau}(q;\mu) \equiv
    \int_q^1
    \frac{\tilde{p}_{t,\tau}(q;\mu)\,\sqrt{b(x;\mu)}}
         {x\,\sqrt{\tilde{p}_{t,\tau}^2(x;\mu) - \tilde{p}_{t,\tau}^2(q;\mu)}}
    \,dx
    = \frac{\Delta}{2}.
    \label{eq:Fermat}
\end{equation}
(The factor $\mu$ cancels between numerator and denominator, confirming that
the angle is finite for $\mu\to0$.)

Writing $\Phi^{t,\tau}(q;\mu) = F_0(q) + \mu\,\Phi_1^{t,\tau}(q) + O(\mu^2)$,
the zeroth-order integral is
\begin{equation}
    F_0(q) = \int_q^1
    \frac{P_0(q)}{x\,\sqrt{P_0^2(x)-P_0^2(q)}}\,dx,
    \label{eq:F0}
\end{equation}
which is the same for both modes (Newtonian limit).
Expanding the integrand in Eq.~\eqref{eq:Fermat} to first order using
Eqs.~\eqref{eq:pt_expand}--\eqref{eq:sqrtb}, and defining
\begin{equation}
    I_0(x;q) = \frac{P_0(q)}{x\,\sqrt{P_0^2(x)-P_0^2(q)}}, \qquad
    \Lambda(x;q) = \frac{P_0^2(x)\,W(x) - P_0^2(q)\,W(q)}{P_0^2(x)-P_0^2(q)},
\end{equation}
one obtains the individual first-order functionals
\begin{align}
    \Phi_1^t(q) &= \int_q^1 I_0(x;q)
        \Bigl[\tfrac{3}{4}W(q) + \tfrac{\bar{m}(x)}{2x} - \tfrac{3}{4}\Lambda(x;q)\Bigr]\,dx,
    \label{eq:Phi1t}\\
    \Phi_1^\tau(q) &= \int_q^1 I_0(x;q)
        \Bigl[-\tfrac{1}{4}W(q) + \tfrac{\bar{m}(x)}{2x} + \tfrac{1}{4}\Lambda(x;q)\Bigr]\,dx.
    \label{eq:Phi1tau}
\end{align}

Note that the $\bar{m}(x)/(2x)$ term (originating from $\sqrt{b}$) is
\emph{identical} in both integrands.
Defining $G[q] \equiv \Phi_1^t(q) - \Phi_1^\tau(q)$ and
$M[q] \equiv \int_q^1 I_0(x;q)\,\bar{m}(x)/(2x)\,dx$, the two
functionals decompose as
\begin{equation}
    \Phi_1^t(q) = \tfrac{3}{4}\,G[q] + M[q],
    \qquad
    \Phi_1^\tau(q) = -\tfrac{1}{4}\,G[q] + M[q],
    \label{eq:Phi1_decomp}
\end{equation}
where
\begin{equation}
    G[q] = \int_q^1 I_0(x;q)\,
    \frac{P_0^2(x)\bigl(W(q)-W(x)\bigr)}{P_0^2(x)-P_0^2(q)}\,dx.
    \label{eq:G_functional}
\end{equation}
$M$ cancels in the difference $G = \Phi_1^t - \Phi_1^\tau$ but enters
the second-order correction below.

\subsection{Perturbative solution to order $\mu^2$}

Write $q^{t,\tau} = q_0 + \mu q_1^{t,\tau} + \mu^2 q_2^{t,\tau} + O(\mu^3)$,
where $q_0$ satisfies $F_0(q_0) = \Delta/2$.
Inserting into Eq.~\eqref{eq:Fermat} and collecting by powers of $\mu$:

\paragraph{Order $\mu$.}
\begin{equation}
    q_1^{t,\tau}\,F_0'(q_0) + \Phi_1^{t,\tau}(q_0) = 0
    \;\Longrightarrow\;
    q_1^{t,\tau} = -\frac{\Phi_1^{t,\tau}(q_0)}{F_0'(q_0)}.
\end{equation}
Using Eq.~\eqref{eq:Phi1_decomp}, the first-order splitting is
\begin{equation}
    \delta q^{(1)} \equiv \mu(q_1^t - q_1^\tau)
    = -\frac{\mu\,G[q_0]}{F_0'(q_0)}.
    \label{eq:dq1}
\end{equation}
For the Schwarzschild interior ($\bar{m}(x)=x^3$, $W(x)=(1-x^2)/2$) this
reduces to the known analytical result
$\delta q^{(1)}\big|_\mathrm{Schw} = \mu q_0(1-q_0^2)^2/4$.

\paragraph{Order $\mu^2$.}
\begin{equation}
    q_2^{t,\tau}\,F_0'(q_0)
    + \frac{(q_1^{t,\tau})^2}{2}\,F_0''(q_0)
    + q_1^{t,\tau}\,(\Phi_1^{t,\tau})'(q_0)
    + \Phi_2^{t,\tau}(q_0) = 0,
    \label{eq:q2_eq}
\end{equation}
where primes on $F_0$ denote $d/dq_0$ and $\Phi_2^{t,\tau}$ are the
second-order angle functionals (corrections to the lapse from the
second-order TOV equation).
We drop $\Phi_2^t - \Phi_2^\tau$ because it requires EOS-specific pressure
integrals at $O(\mu^2)$ and is numerically negligible for $\mu \lesssim 0.3$.
Subtracting the equations for $t$ and $\tau$ and inserting the
decomposition~\eqref{eq:Phi1_decomp}:

\begin{equation}
    \delta q^{(2)} \equiv \mu^2(q_2^t - q_2^\tau) = \mu^2
    \left[
        \frac{G G'/2 + G M' + M G'}{F_0'^{\,2}}
        -\frac{G\bigl(G/2 + 2M\bigr)\,F_0''}{2\,F_0'^{\,3}}
    \right]_{q_0},
    \label{eq:dq2}
\end{equation}
where $G' = dG/dq_0$, $M' = dM/dq_0$ and $F_0'' = d^2F_0/dq_0^2$ are
evaluated at $q_0$ by numerical central-difference derivatives.

\subsection{Summary and implementation}

The complete $O(\mu^2)$ formula is
\begin{equation}
    \boxed{
    \delta q = q^{(t)} - q^{(\tau)}
    = -\frac{\mu\,G[q_0]}{F_0'(q_0)}
    + \mu^2
      \left[
          \frac{GG'/2 + GM' + MG'}{F_0'^{\,2}}
          - \frac{G(G/2+2M)\,F_0''}{2\,F_0'^{\,3}}
      \right]_{q_0}
      + O(\mu^3).
    }
    \label{eq:dq_nlo}
\end{equation}
All functionals ($F_0$, $G$, $M$) are computed from the actual mass profile
$\bar{m}(x)$ of each TOV solution via numerical quadrature; the $q$-derivatives
are evaluated by central finite differences with step $h = 3\times10^{-4}$.
For the analytical Schwarzschild curve, the uniform-density profiles
$\bar{m}(x) = x^3$ and $W(x) = (1-x^2)/2$ are substituted directly.

\paragraph{Relation between terms.}
The second-order correction~\eqref{eq:dq2} contains two contributions:
\begin{enumerate}
    \item \emph{Kinematic feedback}: the term $G G'/2\,/F_0'^2$ arises because
          $q_1^t \neq q_1^\tau$, so their squares differ at $O(\mu^2)$ through
          $F_0''$.
    \item \emph{Profile coupling}: the terms involving $M$ and $G'$ arise because
          the radial-metric correction $\sqrt{b} \approx 1 + \mu\bar{m}/(2x)$
          shifts $\Phi_1^t$ and $\Phi_1^\tau$ by the same amount $M$, but its
          $q$-derivative then appears differently in $q_2^t$ and $q_2^\tau$.
\end{enumerate}

\end{document}
